\documentclass[25pt]{tikzposter}
\usepackage{pifont}
\usepackage[backend=biber,style=vancouver]{biblatex}
\usepackage{xcolor}
\usepackage{booktabs}
\usepackage{tabularx}
\usepackage{enumitem}
\usepackage{graphicx}
\usepackage{multicol}
\usepackage{wrapfig}
\addbibresource{./biblio.bib}
\renewcommand\labelitemi{\ding{118}}
\renewcommand\labelitemii{\ding{239}}

\geometry{paperwidth=36in,paperheight=48in}
\makeatletter
\setlength{\TP@visibletextwidth}{\textwidth-2\TP@innermargin}
\setlength{\TP@visibletextheight}{\textheight-2\TP@innermargin}

\makeatother
\def\posterfaculty{science}
\def\posterbleed{3mm}
\include{uottawa_poster}

\AtBeginDocument{%
    \begin{pgfonlayer}{backgroundlayer}
        \node[
            inner sep=4pt, anchor=south east, fill=white, draw=none,
            rounded corners=5, fill opacity=0, text opacity=1
        ] at (0.5\textwidth-18pt, -0.5\textheight+18pt)%18pt for bleed 3mm, 7 for no bleed
	{\huge{\url{kzhao097@uottawa.ca}}};  %%% ADAPT FOOTNOTE HERE
    \end{pgfonlayer}
}

\newcommand{\RR}{\mathbb{R}}
\newcommand{\NN}{\mathbb{N}}
\newcommand{\OO}{\mathcal{O}}

\newcommand{\MM}{{\rm M}}
\newcommand{\FF}{{\rm F}}
\newcommand{\CC}{{\rm C}}
\newcommand{\EE}{{\rm E}}
\newcommand{\DD}{{\rm D}}
\newcommand{\MAF}{{\rm MAF}}
\newcommand{\Rp}{{\rm R}}
\newcommand{\Sp}{{\rm S}}
\newcommand{\Vp}{{\rm V}}
\newcommand{\Op}{{\rm O}}

\titlegraphic{\includegraphics[height=10em]{uottawa_VER_WHITE}}
\title{\parbox{\linewidth}{\centering \textbf{TriLLIEM}: Triad Log-Linear modelling of Imprinting, Environmental interactions, and Maternal effects}}
\author{Kevin H.\ Zhao and Kelly Burkett}
\institute{Department of Mathematics and Statistics, University of Ottawa}

\begin{document}
\maketitle
\begin{columns}
	\column{0.3}
	\block{Background}
	{\Large
		\begin{itemize}
			\item The trio-based design is used to estimate disease risk in offspring based on the presence of a minor allele and environmental factors.
			\item The Weinberg group~\cite{UmbaWein00,WeinUmba05,Wein+98} proposed a log-linear model to estimate risk from genetic effects.
			\item These effects include:
				\begin{itemize}
					\item Child and maternal genetic effects
					\item Parent-of-origin effects
					\item Gene-environment interactions
				\end{itemize}
			\item The model depends on assumptions about the distribution of mother-father mating type pairs in the population.
			\item No official software for their model was ever released.
		\end{itemize}
	}
	\block{Motivation \& objectives}
	{\Large
		\begin{tikzpicture}[remember picture, overlay]
			\node[anchor=north east]
				at (1.05\linewidth, 0.0\linewidth)
				{\begin{minipage}{0.28\linewidth}
					\centering
				\includegraphics[width=\linewidth]{plots/qr-code}

				\vspace{-0.3em}
				\small\itshape Visit our GitHub repo!
			\end{minipage}
			};
	\end{tikzpicture} \vspace{-0.5em}
		\begin{itemize}
			\item \begin{minipage}[t]{0.7\linewidth} Existing software packages for trio-based designs have limited features, such as requiring a stratified approach to gene-environment interactions.\end{minipage} \\
			\item We developed a new \textsf{R} package, \textbf{TriLLIEM}, which makes full use of the log-linear model.
			\item We used \textbf{TriLLIEM} to conduct a simulation study assessing the model's ability to identify these effects.
			\item We also used \textbf{TriLLIEM} to assess the model's performance when data is sampled from a simulated population with hidden substructures.
		\end{itemize}
	}
	\block{Notation}
	{\large
		\renewcommand{\arraystretch}{1.2}
		\centering
		\begin{tabular}{p{1.8in}>{\raggedleft\arraybackslash}p{7.5in}}
	\toprule
	Genetic factor & Description \\
	\midrule
	$\CC$ & Risk from each copy of child's allele \\
	$\MM$ & Risk from each copy of mother's allele \\
	$I_{\MM}$ & Risk from child inheriting allele from mother \\
	$I_{\FF}$ & Risk from child inheriting allele from father \\
	$\EE:\MM$ & Risk from each copy of mother's allele conditional on environmental exposure \\
	$\EE:I_{\MM}$ & Risk from child inheriting allele from mother conditional on environmental exposure \\
	$\EE:I_{\FF}$ & Risk from child inheriting allele from father conditional on environmental exposure \\
	\bottomrule
\end{tabular}

\vspace{2.35em}

\begin{tabular}{p{1.8in}>{\raggedleft\arraybackslash}p{7.5in}} 
	\toprule 
	Mating type model & Description \\ 
	\midrule 
	HWE & Hardy-Weinberg equilibrium mating type model, parent genotypes are independent and a result of random draw of alleles from population. \\
	MS & Mating symmetry mating type model, assumes mother-father mating type probabilities are interchangeable. \\ 
	MaS & Mating asymmetry mating type model, probabilities of each mother-father mating type pair are independent. \\
	\bottomrule
	\end{tabular}
	}

	\column{0.7}
	\block{Simulation study 1 --- Assessing effect estimation and power}
	{\LARGE
	\begin{itemize}
		\item Under various true genetic effects, we simulated and analyzed $2,000$ data sets with $1,000$ case-triads and $1,000$ control-triads, and another $2,000$ data sets with only $2,000$ case-triads.
		\item With $50\%$ case- and control-triads in a population under mating asymmetry, we observed high Type 1 error rates for child effects when modelling maternal imprinting by environment interactions compared to using the same sample size with $100\%$ case-triads (Left figure, top row).
		\item The non-stratified (unhatched) approach that \textbf{TriLLIEM} permits has substantially greater power than the stratified (hatched) approach other packages use for gene-environment interactions (Right figure).
	\end{itemize}\vspace{0.5em}
	\begin{subcolumns}
		\subcolumn{0.5}
		\centering
		\includegraphics[width=0.5\linewidth]{plots/reject_mas_mimpE.pdf}
		\subcolumn{0.5}
		\includegraphics[width=0.5\linewidth]{plots/stratcomp_Im.pdf}
		\vspace{-3.405em}
	\end{subcolumns}
	}
	\block{Simulation study 2 --- Effect of population stratification}
	{\Large
		\begin{minipage}{0.525\linewidth}
\begin{tikzfigure}
	\begin{tabular}{p{2in}p{2.25in}p{3in}p{2in}p{2in}}
  \hline
  Sub-population       & Minor allele frequency   & Environmental exposure rate & Disease prevalence & Population proportion \\ \hline
$1$  & $20\%$ & $30\%$        & $10\%$    & $40\%$    \\
$2$ & $40\%$ & $50\%$        & $20\%$    & $60\%$    \\ \hline
\end{tabular}

\vspace{0.5em}

\end{tikzfigure}
		\begin{itemize}
		\item The population consisted of two hidden sub-populations with parameters given in the table above.
		\item Simulations used one data set with $1,000$ case-triads and $1,000$ control-triads, and another with only $2,000$ case-triads.
		\item The log-linear model was run on each data set with a variety of genetic effects in the model.
		\item This was repeated $2,000$ times, and the average Type 1 error rates for the model parameters are given in the tables below for simulations with both case- and control-triads (top) and with only case-triads (bottom).
		\end{itemize}
	\end{minipage}
	\begin{minipage}{0.45\linewidth}
\centering
\resizebox{\ifdim\width>\linewidth\linewidth\else\width\fi}{!}{
\begin{tabular}[t]{llllllll}
\toprule
$\CC$ & $\MM$ & $\EE:\MM$ & $\EE:I_{\MM}$ & $\EE:I_{\FF}$ & $I_{\MM}$ & $I_{\FF}$ & Mating type model\\
\midrule
0.2775 & 0.2345 & 0.0670 &  &  &  &  & HWE\\
0.4825 & 0.4675 &  & 0.0670 &  & 0.1915 &  & HWE\\
0.0570 & 0.4670 &  &  & 0.055 &  & 0.257 & HWE\\
\addlinespace
0.2625 & 0.2175 & 0.0630 &  &  &  &  & MS\\
0.4485 & 0.4285 &  & 0.0665 &  & 0.1835 &  & MS\\
0.0660 & 0.4295 &  &  & 0.057 &  & 0.242 & MS\\
\addlinespace
0.1650 & 0.3035 & 0.0645 &  &  &  &  & MaS\\
0.2945 & 0.4805 &  & 0.0610 &  & 0.1325 &  & MaS\\
0.0600 & 0.4795 &  &  & 0.066 &  & 0.176 & MaS\\
\bottomrule
\end{tabular}}


\vspace{0.4em}

\centering
\resizebox{\ifdim\width>\linewidth\linewidth\else\width\fi}{!}{
\begin{tabular}[t]{llllllll}
\toprule
$\CC$ & $\MM$ & $\EE:\MM$ & $\EE:I_{\MM}$ & $\EE:I_{\FF}$ & $I_{\MM}$ & $I_{\FF}$ & Mating type model\\
\midrule
0.0510 & 0.0475 & 0.052 &  &  &  &  & HWE\\
0.0435 & 0.0450 &  & 0.0505 &  & 0.0500 &  & HWE\\
0.0485 & 0.0450 &  &  & 0.0425 &  & 0.0445 & HWE\\
\addlinespace
0.0540 & 0.0510 & 0.056 &  &  &  &  & MS\\
0.0510 & 0.0505 &  & 0.0565 &  & 0.0555 &  & MS\\
0.0565 & 0.0500 &  &  & 0.0475 &  & 0.0530 & MS\\
\bottomrule
\end{tabular}}

\end{minipage}
	\vspace{-0.5em}
}

	\block{Discussion}
	{\LARGE
	\begin{itemize}
		\item Our package, \textbf{TriLLIEM}, is a complete implementation of the log-linear model, allowing researchers to no longer be restricted to the stratified approach for gene-environment interactions.
		\item We verified that in populations with hidden sub-populations, the model performs well when given data with only case-triads, and only suffers small inflation in identifying gene-environment interactions when control-triads are included.
	\end{itemize}
	\vspace{-0.5em}
}
	\block{References}{
	\printbibliography[heading=none]
}

\end{columns}
\end{document}
